\section{Additional experimental details}
\label{app:repro}

\paragraph{Pretraining and linear evaluation.}
The input crop is $256\times256$, with patch size 16 and a ViT-Base encoder.
The predictor has depth 6 and hidden dimension 384.
Pretraining uses AdamW with peak learning rate $0.00025$, a five-epoch
warmup, final learning rate $0.000001$, weight decay scheduled from $0.04$
to $0.4$, and EMA coefficient from $0.996$ to $1.0$.
The effective batch size is 512, formed from batches of 64 with gradient
accumulation. Rectangle target scales are $0.15$--$0.20$ of image area,
aspect ratios $0.75$--$1.5$, and context scales $0.85$--$1.0$.
Random resized crops use area scale $0.3$--$1.0$.
The guidance probability increases over the epoch-25 to epoch-30
continuation window; fixed-length anatomy targets use $K=16$ throughout.
The frozen linear head uses AdamW, learning rate $0.0004$, weight decay
$0.05$, batch size 256, seed 42, and a 50-epoch budget with validation-AUC
selection and patience 15.

\paragraph{Guide sources.}
\textsc{envelope} uses the MIRAGE-derived repaired hard retinal envelope,
produced by MIRAGE-Large (frozen encoder, ConvNeXt decoder) fine-tuned only on
GOALS layer annotations (background, nerve fibre layer, ganglion cell and inner
plexiform layers, choroid). The anatomy-shaped and coverage strategies use
precomputed soft guides from
MIRAGE with frozen residual adapters. The anatomy-v1 guide uses a
decoder-feature adapter; anatomy-v2 and \textsc{cover} use an encoder-feature
adapter. These guide products are not independently controlled in the
cross-strategy comparison. \ArmBest{} uses image intensity rather than a
segmentation model: in the trained configuration its band
(Equation~\ref{eq:centroid}) is $\CRBandRows$ rows high and spans the central
$\CRBandCols$ columns of the patch grid (region fraction $\CRBandRegion$,
lateral fraction $\CRBandLateral$).
The adapters optimize feature-similarity objectives on FairVision-derived
image caches, not a glaucoma classification loss, and are frozen during
I-JEPA training. The retained provenance does not independently establish
subject-level non-overlap for all data used to train the guidance models.

\paragraph{Included and excluded continuations.}
\label{app:excluded}
The early anatomy-v1 configuration resumes an \textsc{envelope} epoch-27
checkpoint rather than a documented epoch-25 fork, and its producing launch record
is incomplete. Its only result is at epoch 30, where it lies
$\DAnatomyOneEnvelopeEpThirty$ above \textsc{envelope} (CI
$\DAnatomyOneEnvelopeEpThirtyCI$). This comparison inherits the lineage
uncertainty.
Anatomy-v2 results use the continuation with full-precision target-encoder
outputs. Probes from its earlier mixed-target-precision continuation are
excluded, as are earlier uniform controls that used different target
precision and context selection. The main table retains the comparable
policy endpoints, including the adverse \textsc{cover} trajectory.

\paragraph{Mask measurement.}
\label{app:geomprov}
Table~\ref{tab:geometry} replays the sampler code in each strategy's trained
configuration (commit {\footnotesize\texttt{\CRSamplerCommit}}), including
the \ArmBest{} band above and \textsc{envelope} placement on the hard envelope
with the fallback rule used in training. The scope is $\CRAuditVolumes$
training volumes with $\CRAuditSlices$ views each; every strategy sees the same
crops, block sizes and placement seeds. Tissue is guide occupancy of at least
$0.25$ on the soft MIRAGE envelope, on the $16\times16$ patch grid. The
submitted tables used $\CROldSlices$ slices from $\CROldVolumes$ volumes with
seed $\CROldSeed$, sizes drawn separately for each strategy, the code-default
\ArmBest{} band width and \textsc{envelope} placed on the soft guide.
Table~\ref{tab:delivered} compares both. The band width has little effect on
delivered context; most of the difference comes from replay noise and
pairing, and for \textsc{envelope} also from the guide product and fallback
rule.

\begin{table}[h]
\centering
\caption{Submitted versus camera-ready mask geometry. Each cell gives the
submitted value, then the value under the trained configuration and the
protocol of Table~\ref{tab:geometry}. Per-image columns are measured one view
at a time ($\CRImgDraws$ single-view draws per seed), before batching;
the last column is the context delivered at batch size $\CRAuditBatch$.}
\label{tab:delivered}
\footnotesize
\setlength{\tabcolsep}{3pt}
\begin{tabular}{lcccccc}
\toprule
& \multicolumn{5}{c}{per image} & delivered \\
\cmidrule(lr){2-6}
policy & tissue hidden & purity & mask ratio & context & loss slots & context \\
\midrule
\textsc{random}     & \CRImgHidOldRandom{} / \CRImgHidRandom         & \CRImgPurOldRandom{} / \CRImgPurRandom         & \CRImgMaskOldRandom{} / \CRImgMaskRandom         & \CRImgCtxOldRandom{} / \CRImgCtxRandom         & \CRImgSlotsOldRandom{} / \CRImgSlotsRandom         & \CRCtxOldRandom{} / \CRCtxRandom \\
\ArmBest{}          & \CRImgHidOldCentroid{} / \CRImgHidCentroid     & \CRImgPurOldCentroid{} / \CRImgPurCentroid     & \CRImgMaskOldCentroid{} / \CRImgMaskCentroid     & \CRImgCtxOldCentroid{} / \CRImgCtxCentroid     & \CRImgSlotsOldCentroid{} / \CRImgSlotsCentroid     & \CRCtxOldCentroid{} / \CRCtxCentroid \\
\textsc{envelope}   & \CRImgHidOldEnvelope{} / \CRImgHidEnvelope     & \CRImgPurOldEnvelope{} / \CRImgPurEnvelope     & \CRImgMaskOldEnvelope{} / \CRImgMaskEnvelope     & \CRImgCtxOldEnvelope{} / \CRImgCtxEnvelope     & \CRImgSlotsOldEnvelope{} / \CRImgSlotsEnvelope     & \CRCtxOldEnvelope{} / \CRCtxEnvelope \\
\textsc{cover}      & \CRImgHidOldCover{} / \CRImgHidCover           & \CRImgPurOldCover{} / \CRImgPurCover           & \CRImgMaskOldCover{} / \CRImgMaskCover           & \CRImgCtxOldCover{} / \CRImgCtxCover           & \CRImgSlotsOldCover{} / \CRImgSlotsCover           & \CRCtxOldCover{} / \CRCtxCover \\
\textsc{anatomy-v2} & \CRImgHidOldAnatomyTwo{} / \CRImgHidAnatomyTwo & \CRImgPurOldAnatomyTwo{} / \CRImgPurAnatomyTwo & \CRImgMaskOldAnatomyTwo{} / \CRImgMaskAnatomyTwo & \CRImgCtxOldAnatomyTwo{} / \CRImgCtxAnatomyTwo & \CRImgSlotsOldAnatomyTwo{} / \CRImgSlotsAnatomyTwo & \CRCtxOldAnatomyTwo{} / \CRCtxAnatomyTwo \\
\bottomrule
\end{tabular}
\end{table}

\paragraph{Data use and ethics.}
\label{app:ethics}
This is secondary research on the public, de-identified FairVision release
under its CC BY-NC-ND 4.0 data licence. No new participants were recruited
or identifying fields added. A separate institutional ethics or exemption
determination is not recorded for this project. The models are not clinically
validated or intended for patient care. The illustrative OCT image in the
main paper comes from the separately attributed CC BY 4.0 OCTDL publication,
not from the FairVision evaluation cohort.

\section{Registered replication protocol}
\label{app:replication}
The replication plan was committed to the project repository before any new
continuation produced a checkpoint and before any new probe (commit
{\footnotesize\texttt{86ae8823f3afc9fa2c9ed6a3cb0c459b513ba1c5}}).
\textsc{random}, \ArmBest{} and \textsc{envelope} each continue from the shared
epoch-25 checkpoint with two new seeds, which set data order, crops, mask draws
and data-loader streams. Runs with the same seed share data order only.
\textsc{envelope} uses the fallback placement of the original run. A
uniform-placement control receives \ArmBest{}'s per-image budgets with the first
seed. Optimiser, schedule horizon and effective batch size match the original
continuations, with full-precision teacher targets as in the originals; each run
stops after its epoch-50 checkpoint. The original \textsc{random} and \ArmBest{}
continuations ran on different GPU hardware from the original \textsc{envelope}
continuation and all new runs.
The probe keeps the main protocol but runs in fp32, seeds each of five linear
heads after the feature caches are created, and reports the mean AUC over head
seeds together with their s.d.; the original encoders are re-probed in the same
way. Test predictions are stored sealed until the results freeze, and decisions
during the campaign use validation AUC only. For each new seed the primary
contrasts are \ArmBest{} and \textsc{envelope} minus \textsc{random}, with paired
test-case bootstrap intervals. A guided strategy is labelled consistent in
direction only if both new-seed differences are positive and their mean exceeds
the range of the three \textsc{random} realizations. A positive mean otherwise
counts as within run-to-run variation, and a non-positive mean as reversed.
Analyses including the original runs, validation-split differences and pooling
variants are secondary. Every started run is reported.

\section{Guide agreement and envelope plausibility}
\label{app:guides}
FairVision has no retinal layer annotations, so neither guide can be scored for
accuracy on it. We fixed failure rules before computing any rate and compared
the \ArmBest{} band with the hard MIRAGE envelope (patch occupancy at least
$0.25$) on full B-scans of $\CRGuideVolumes$ training volumes,
$\CRGuidePerStratum$ from each combination of glaucoma label and tertile of a
signal-to-noise proxy, with ten B-scans per volume ($\CRGuideBscans$ B-scans).
The computation was repeated on all $\CRCensusVolumes$ training volumes. A
B-scan counts as a gross disagreement when the band--envelope IoU is below
$0.2$ or more than half of the band cells contain no envelope. Rates are
population-weighted, with 95\% volume-cluster bootstrap intervals.

The band disagreed grossly with the envelope in $\CRFOne\%$ of B-scans
(\CRFOneCI, all training volumes $\CRFOneCensus\%$) and was off-centre in at
least three of ten columns in $\CRFTwo\%$. Disagreement was more frequent in
glaucoma than in normal eyes ($\CRFOneGlaucoma\%$ versus $\CRFOneNormal\%$) and
in the lowest than in the highest signal tertile ($\CRFOneLowSNR\%$ versus
$\CRFOneHighSNR\%$). The median IoU was $\CRIoUMedian$. Because the envelope is
typically about half as thick as the $\CRBandRows$-row band, an IoU of this
size is expected even for a centred band. Of the envelopes, $\CREnvAnyFlag\%$
(\CREnvAnyFlagCI) raised at least one plausibility flag: three or more
components ($\CREnvFragmented\%$), irregular thickness ($\CREnvThickness\%$) or
implausible position ($\CREnvPosition\%$). These are agreement and plausibility
statistics between two automatic guides, not anatomical accuracy. The envelope
model's only quantitative validation is in-domain on GOALS (foreground Dice
$\CRMirageDice$).
